\newpage
\chapter{KESIMPULAN DAN SARAN} \label{Bab V}

\section{Kesimpulan} \label{V.Kesimpulan}
Berdasarkan hasil analisis penelitian yang telah dilakukan, dapat dirumuskan kesimpulan sebagai berikut:

\begin{enumerate}
  \item Arsitektur VideoMAE terbukti sangat sangkil untuk tugas klasifikasi gerakan visual Bahasa Isyarat SIBI. Konfigurasi optimal (E12) dicapai melalui kombinasi arsitektur ViT-Base, bobot pra-latih Kinetics-400, \textit{dense sampling}, tanpa augmentasi spasio-temporal campuran, serta strategi \textit{two-stage fine-tuning}. Konfigurasi ini menghasilkan metrik performa tertinggi dengan akurasi rata-rata sebesar 97.18\% dan capaian tertinggi pada satu \textit{fold} sebesar 98.92\%.
  
  \item Penggunaan bobot pra-latih (\textit{pre-trained weights}) merupakan prasyarat mutlak bagi model berbasis Transformer pada tugas visi komputer, khususnya untuk klasifikasi bahasa isyarat. Hal ini dibuktikan secara empiris oleh penurunan performa yang drastis ketika model dilatih dari inisialisasi acak (anjlok dari 63.85\% menjadi kisaran 11--12\%). Arsitektur Transformer pada dasarnya kekurangan \textit{inductive bias} (seperti invariansi translasi spasial) yang secara alami dimiliki oleh \textit{Convolutional Neural Network} (CNN), sehingga model Transformer sangat bergantung pada representasi awal dari tahap pra-latih agar dapat berkonvergensi secara efektif pada kondisi dataset berskala kecil.
  
  \item Penerapan augmentasi data tingkat lanjut yang memanipulasi komposisi piksel (\textit{mixup, CutMix, reprob}) terbukti bersifat destruktif terhadap performa pengenalan, di mana penonaktifan teknik tersebut memicu lonjakan performa sebesar 30.09 poin persentase. Secara teoritis, modifikasi dan fusi fitur lintas-kelas tersebut merusak integritas spasio-temporal dan detail artikulasi gestur (seperti bentuk konfigurasi tangan dan lintasan gerak) yang merupakan elemen diskriminatif esensial dalam struktur linguistik bahasa isyarat.
  
  \item Strategi penarikan sampel temporal (\textit{temporal sampling}) berpengaruh signifikan dalam mengekstraksi dinamika rentang gestur isyarat. Pada komparasi terkontrol, pendekatan \textit{uniform sampling} terbukti lebih reliabel daripada \textit{dense sampling} karena kemampuannya dalam merepresentasikan keseluruhan fase transisi gestur dari awal hingga akhir secara merata. Walau demikian, optimalisasi akhir dari strategi \textit{sampling} bersifat dinamis dan turut ditentukan oleh kapasitas pemrosesan \textit{backbone} serta dataset asal prapelatihannya.
  
  \item Evaluasi lintas paradigma arsitektur mengonfirmasi bahwa pendekatan \textit{self-supervised video transformer} (VideoMAE E12: 97.18\%) berhasil melampaui paradigma \textit{3D CNN} (I3D: 96.86\%) maupun \textit{supervised transformer} (ViViT: 94.81\%). Hal ini memvalidasi keunggulan arsitektur Transformer dalam memodelkan dependensi spasio-temporal jarak jauh pada data berdimensi tinggi, asalkan didukung representasi prapelatihan yang kaya. Di samping itu, isyarat dengan kedekatan komponen visual secara spasial, seperti kelas ``Mau'' dan ``Suka'', secara konsisten teridentifikasi sebagai batas kinerja (\textit{performance bottleneck}) pada seluruh ragam arsitektur, yang mengindikasikan tingginya kebutuhan terhadap mekanisme atensi beresolusi tinggi pada domain bahasa isyarat.
\end{enumerate}

\section{Saran} \label{V.Saran}
Berdasarkan temuan dan keterbatasan yang diperoleh selama penelitian, diajukan beberapa saran untuk pengembangan selanjutnya:

\begin{enumerate}
  \item Perluasan dataset. Dataset SIBI yang digunakan dalam penelitian ini terdiri dari 924 klip video untuk 10 kelas kata dasar. Untuk meningkatkan generalitas model, disarankan memperluas cakupan kosakata SIBI ke setidaknya 50--100 kata dengan melibatkan lebih banyak responden dari berbagai latar belakang usia, jenis kelamin, dan gaya isyarat. Keragaman responden akan membantu model belajar fitur yang lebih invarian terhadap variasi antar-penutur.

  \item Eksplorasi augmentasi yang tepat. Temuan bahwa \textit{mixup}, \textit{CutMix}, dan \textit{reprob} bersifat kontraproduktif membuka peluang untuk menyelidiki augmentasi alternatif yang tidak merusak integritas spasio-temporal gestur. Augmentasi yang beroperasi murni pada domain tampilan (\textit{appearance-only}) dinilai lebih aman untuk data bahasa isyarat, misalnya \textit{brightness} dan \textit{contrast jitter} untuk mensimulasikan variasi pencahayaan, \textit{random spatial crop} untuk meningkatkan ketahanan terhadap perbedaan framing dan posisi, serta penambahan \textit{Gaussian noise} atau \textit{blur} ringan untuk mensimulasikan variasi kualitas kamera. Ketiga teknik tersebut memanipulasi atribut piksel tanpa mengubah urutan atau struktur temporal gestur yang menjadi informasi kritis pada bahasa isyarat.

  \item Pengujian pada kondisi tidak terkontrol. Seluruh eksperimen dalam penelitian ini dilakukan pada dataset yang direkam dalam kondisi terkontrol (pencahayaan stabil, latar polos). Disarankan menguji kinerja model pada video yang direkam dalam kondisi dunia nyata dengan variasi pencahayaan, latar belakang kompleks, dan sudut pandang kamera yang berbeda untuk mengukur ketahanan sistem secara lebih realistis.

  \item Eksplorasi strategi \textit{fine-tuning} lanjutan. Penelitian ini menggunakan \textit{two-stage fine-tuning} dengan konfigurasi \textit{learning rate} dan epoch yang tetap. Teknik lanjutan seperti \textit{layer-wise learning rate decay} atau \textit{progressive unfreezing} berpotensi meningkatkan kinerja lebih lanjut tanpa menambah data berlabel, karena keduanya memungkinkan lapisan awal backbone tetap stabil sementara lapisan atas diadaptasi secara bertahap terhadap domain bahasa isyarat.

  \item Eksplorasi \textit{positional encoding} untuk kelas dengan kemiripan visual tinggi. Hasil eksperimen menunjukkan bahwa kelas Mau dan Suka secara konsisten menjadi sumber kesalahan klasifikasi terbesar pada seluruh model yang diuji. Kedua isyarat tersebut memiliki komponen gerakan tangan yang secara visual serupa, sehingga model perlu mengandalkan perbedaan subtil dalam urutan temporal dan posisi gestur. Disarankan untuk mengeksplorasi teknik \textit{relative positional encoding} atau \textit{learnable temporal embedding} yang memungkinkan model menangkap relasi urutan antar-frame secara lebih eksplisit, sehingga perbedaan sekuensial antar isyarat yang secara semantik berdekatan dapat direpresentasikan dengan lebih diskriminatif.

  \item Pengembangan sistem berbasis \textit{real-time}. Penelitian ini difokuskan pada tahap analisis kinerja klasifikasi. Sebagai langkah lanjutan, arsitektur terbaik (VideoMAE E12 atau I3D) dapat dikembangkan menjadi sistem pengenalan bahasa isyarat berbasis \textit{real-time} yang terintegrasi dengan antarmuka pengguna, sehingga temuan penelitian ini dapat memberikan dampak langsung bagi komunitas Tuli di Indonesia.
\end{enumerate}
